\chapter{Epilogue}
\label{chapter:type_theory}
\label{ch:typetheory}

\section{Conclusion}

Originally, this thesis aimed to give a complete categorical account of e-graphs, including sharing with the behaviour exhibited by traditional e-graphs, as well as cycles.
It was initially expected that $\catname{SLat}$-enrichment preserved cartesian structure and that, moreover, the product in an $\catname{SLat}$-category would be given using the tensor product of semilattices.
This expectation was partly motivated by the fact that the functor $\widetilde{F}_{\catname{SLat}}$ lifts an ordinary adjunction to a $\catname{SLat}$-adjunction. 
Although this statement is correct, the particular $\catname{SLat}$-functors participating in the resulting adjunction are important for determining the induced categorical structure.
Thus, as we established in~\cref{prop:binary_products_not_preserved}, the product is not preserved by the free enrichment and, furthermore, the product has the wrong property for e-graph semantics.

The treatment under the original assumption allowed for particularly nice string diagrams that matched the traditional e-graph almost 1-to-1.
For example, one could draw a diagram for traced cartesian $\catname{SLat}$-categories of the form:
\[
\adjustbox{max width=0.8\textwidth}{
  \tikzfig{./figures/epilogue/e-graph-example-2}
}
\]
where such a string diagram is shown on the right along with the corresponding e-graph on the left. 
Note that the e-graph is cyclic and, in particular, represents terms $a$, $a * 1$, $(a * 1) * 1, \ldots$. 
This is not the case for the string diagram.
To illustrate this, consider a simpler diagram (1) below:
\[
\adjustbox{max width=0.8\textwidth}{
  \tikzfig{./figures/epilogue/trace}
}~,
\]
which would intuitively have corresponded to the e-graph that represents all terms $a, f(a), f(f(a))$ and so on.
This is not the case in a cartesian traced $\catname{SLat}$-category.
First, we apply the naturality of the copy $\comonoid$ to get the diagram (2).
Then we apply the product structure of the $\catname{SLat}$-category (\emph{cf.}~\cref{prop:binary_products_not_preserved}) and get the diagram (3).
The diagram (3) represents only terms $a$ and $f(f(f(f(f(\ldots)))))$.

Consequently, the categorical construction developed in this thesis captures a fragment of traditional e-graphs rather than their full cyclic and sharing structure.
This fragment is still interesting and was successfully generalised to the monoidal setting.
Sharing in this fragment is represented by a non-natural comonoid such that, if the component is an e-box, all references to a shared element are resolved to the same choice, which might be a reasonable semantics for some variation of e-graphs.

We characterised the generalised e-graphs as morphisms in free semilattice-enriched symmetric monoidal categories and gave them a concrete combinatorial representation as extended cospans of e-hypergraphs.
We proved that this representation is sound and complete and established a correspondence between enriched term rewriting and extended double-pushout rewriting with interfaces, modulo the relevant structural rules.
This provided a graph-rewriting account of equality saturation for symmetric monoidal theories.
We also extended this framework to handle variable binding by introducing closed e-hypergraphs, which were also shown to be sound and complete with respect to the class of terms they represent.
We gave a lot of examples comparing the traditional e-graphs and the generalised e-graphs.

Future work could consider supporting the right notion of sharing and cycles.
A possible direction is to investigate categorical accounts of iteration and recursive definitions, including iteration theories~\cite{bloomIterationTheories1993,goncharovShadesIterationElgot2023}, Elgot theories~\cite{adamekElgotTheoriesNew2011}, and Elgot monads~\cite{goncharovCompleteElgotMonads2016}.
Another interesting direction would be to investigate the algorithmic aspects of EDPOI rewriting modulo $\catname{SLat}$-structure.

\section{Type theory for \texorpdfstring{$\catname{SLat}$}{SLat}-multicategories}
\label{sec:type-theory-epilogue}

This section presents an idea that emerged from the development in the
previous chapters: e-graphs may be approached type-theoretically, using the
bridge between category theory and type theory.  In~\cref{chapter:algebras} we
observed that the terms represented by an e-graph come from an algebraic
theory, and that algebraic theories may be viewed as categories with finite
products.  For the purpose of giving syntax to their morphisms, however,
multicategories are more natural.  They take a morphism with several inputs
and one output as primitive, and therefore let us write $f(x_1,\ldots,x_n)$
without first encoding the tuple of arguments using tensors, symmetries,
identities, and composition.

The enriched category theory of~\cref{chapter:enriched} extends to this
multicategorical setting.  In particular, a morphism in an
$\catname{SLat}$-multicategory may be joined with another morphism of the same
profile, and multicategorical composition distributes over this join.  This
is precisely the behaviour wanted of an e-class: $M\join N$ records two
alternative representatives, while applying a constructor to e-classes
forms all combinations of their representatives.

A type theory for $\catname{SLat}$-multicategories makes this interpretation
usable.  Its terms are considerably easier to read than raw monoidal
$\Sigma$-terms, and its inference rules can absorb some of the categorical
equations.  Identity substitution, for example, becomes literal
substitution of a variable rather than an equation involving an identity
box.  Such a language could eventually provide a basis for reasoning about
equality saturation type-theoretically, or for a programming language with
first-class collections of equivalent terms.

This was the goal of an earlier development undertaken during this thesis.
It attempted to model ordinary e-graphs using cartesian
$\catname{SLat}$-multicategories, with explicit \texttt{let} bindings for
sharing and \texttt{letrec} for cycles.  A problem with that approach was
discovered too late for the development to become a thesis contribution:
free $\catname{SLat}$-enrichment does not preserve cartesian structure
(\cref{prop:binary_products_not_preserved}).  More importantly, contraction
is incompatible with the intended reading of join and ordinary textual
substitution.  The full cartesian and recursive calculi are therefore not
claimed here.

There is nevertheless a coherent fragment left behind by that development.
Free enrichment does preserve \emph{affine} structure---exchange and
weakening, but not contraction.  We sketch below a type theory for
representable affine $\catname{SLat}$-multicategories, the coherence and
substitution arguments suggested by its Agda formalisation, and the
initiality theorem that we expect it to satisfy.  Restricting further from
affine maps to bijections gives a symmetric calculus whose terms may be used
as a conventional syntax for the string diagrams encountered in
\cref{chapter:monoidal}.

\subsection{Why categorical type theories?}

The familiar example of a categorical type theory is the simply typed
$\lambda$-calculus, whose categorical semantics is given by cartesian closed
categories.  The general lesson is not confined to that example: categorical
structure may be presented by judgments and inference rules, with types as
objects, contexts as input interfaces, and derivations as morphisms.
Composition is then implemented by substitution.  Proving that substitution
is unital and associative proves the corresponding categorical axioms once,
at the level of syntax.

We follow here the point of view developed systematically by Shulman in
\emph{Categorical Logic from a Categorical Point of View}
\cite{shulman2016categorical}: rather than first writing an unrestricted term
language and then imposing every structural equation, one chooses the shape
of contexts and inference rules so that the desired categorical structure is
built into derivations.  This is useful for e-graphs for two reasons.  First,
it replaces verbose monoidal expressions by ordinary many-input terms.
Secondly, enrichment becomes a term former.  The semilattice operation is
written directly as $M\join N$, so congruence is the multilinearity of the
inference rules rather than a separate family of rules for every operation
symbol.

There is also a conceptual benefit.  Open and closed e-graphs can be treated
uniformly: a judgment
\[
  x_1:A_1,\ldots,x_n:A_n\vdash M:B
\]
represents an e-class with an interface.  If the calculus has an initial
categorical semantics, interpretations of this syntax in other models are
then obtained from its universal property rather than defined separately for
each language.

\subsection{The multicategorical setting}

Multicategories were introduced by Lambek~\cite{LambekMulticategories} as a
generalisation of categories in which a morphism may have a finite list of
inputs and one output.  Thus, for objects $A_1,\ldots,A_n,B$, a multicategory
$\mathcal M$ has a set
\[
  \mathcal M(A_1,\ldots,A_n;B)
\]
of multimorphisms.  An ordinary category is the unary part of this structure.
Multicategorical composition substitutes a multimorphism into an input of
another one; its laws are precisely the unit, associativity, and interchange
laws expected of capture-avoiding substitution.  The full definition and the
placed-composition equations are recalled in
Appendix~\ref{appendix:type-theory-multicategories}.

A multicategory is \emph{representable} when every finite list of inputs has
a tensor product representing it.  In the binary case this means that there
is an object $A\otimes B$ and a universal multimorphism
\[
  \xi_{A,B}\in\mathcal M(A,B;A\otimes B)
\]
such that composition with $\xi_{A,B}$ gives natural bijections
\[
 \mathcal M(\Gamma,A\otimes B,\Delta;C)
 \cong
 \mathcal M(\Gamma,A,B,\Delta;C).
\]
There is similarly a representing unit $\I$ for the empty list.  By a theorem
of Hermida~\cite[Theorem~9.8]{hermida_representable_2000}, representable
multicategories and monoidal categories carry the same information, up to the
appropriate $2$-categorical equivalence.  The advantage of the
multicategorical presentation is syntactic: several named inputs remain
visible instead of being hidden inside a tensor product.

Additional structural rules are described by actions on the input list.  An
action of bijections permits exchange and corresponds, in the representable
case, to symmetric monoidal structure.  An action of all functions also
permits weakening and contraction and corresponds to cartesian structure.
For the surviving development we use only injections.  They permit inputs to
be renamed, permuted, and left unused, but never identify two inputs.  We call
the resulting structure \emph{affine}; representably, it corresponds to a
symmetric monoidal category whose tensor unit is terminal.  Enriching this
structure in $\catname{SLat}$ makes every multimorphism set a semilattice and
requires composition and the structural actions to preserve joins.  Concrete
definitions are given in
Appendix~\ref{appendix:type-theory-affine-multicategories}; see also Johnson and
Yau~\cite{johnson2023homotopytheoryenrichedmackey} for enriched
multicategories in general.

\subsection{Why only the affine fragment survives}

The obstruction to cartesian structure can be seen without any syntax.  If
$\mathcal C$ has a product $A\times B$, free semilattice enrichment replaces a
hom-set by the semilattice $\mathcal P_f^+$ of its finite non-empty subsets.
Consequently,
\[
 \mathcal P_f^+\bigl(\mathcal C(X,A)\times\mathcal C(X,B)\bigr)
 \not\cong
 \mathcal P_f^+\bigl(\mathcal C(X,A)\bigr)
 \times
 \mathcal P_f^+\bigl(\mathcal C(X,B)\bigr)
\]
in general.  The map from left to right remembers only the two projections
and forgets which choices were correlated.  For example, the relations
\[
 \{(f,g),(f',g')\}
 \quad\text{and}\quad
 \{(f,g'),(f',g)\}
\]
have the same projections but need not be equal.

The same problem appears syntactically.  Were contraction allowed, substituting
$M\join N$ for both occurrences of $x$ in $f(x,x)$ would produce
\[
 f(M,M)\join f(M,N)\join f(N,M)\join f(N,N),
\]
whereas copying the already chosen alternative would produce only
$f(M,M)\join f(N,N)$.  Thus contraction confuses an independent choice with
a correlated one.  Weakening causes no such mixed terms: discarding
$M\join N$ simply gives the same result twice, which idempotence identifies.
This is the exact boundary between the abandoned cartesian development and
the affine calculus below.  In particular, the calculus is informative about
the structure of e-graphs, but it is not by itself a syntax for their general
sharing behaviour.

\subsection{The calculus}

Let $\mathcal G$ be a multigraph: it has a set of objects and a set
$\mathcal G(A_1,\ldots,A_n;B)$ of generating arrows for each input profile and
output.  We regard it as a many-sorted signature.  Types and contexts are
generated by
\[
 A ::= a\;(a\in\obj{\mathcal G})\mid\I\mid A\otimes B,
 \qquad
 \Gamma ::= ()\mid\Gamma,x:A,
\]
where the variables in a context are distinct.  Contexts are considered up
to injective renaming, so their displayed order is merely a choice of
enumeration.

We write $\sum_{i=1}^{m}M_i$ for the non-empty iterated join
$M_1\join\cdots\join M_m$; the symbol $\sum$ does not introduce an additional
term former.  Each displayed $M_i$ is join-free.  Constructors are extended
over joins by taking all combinations of their summands.  The derivation rules
are the following; all contexts juxtaposed in one conclusion are disjoint,
and bound variables are taken up to $\alpha$-renaming.
\begin{gather*}
 \inferrule*[right=$\mathsf{id}$]{\;}{x:A\vdash x:A}
 \qquad
 \inferrule*[right=$\mathsf{wk}$]
 {\Gamma\vdash\sum_iM_i:A}
 {\Gamma,x:B\vdash\sum_iM_i:A}
 \\
 \inferrule*[right=$f$]
 {f\in\mathcal G(A_1,\ldots,A_n;B)\quad
  \Gamma_l\vdash\sum_{j=1}^{m_l}M_{lj}:A_l\;(1\leq l\leq n)}
 {\Gamma_1,\ldots,\Gamma_n\vdash
  \displaystyle\sum_{j_1=1}^{m_1}\cdots\sum_{j_n=1}^{m_n}
  f(M_{1j_1},\ldots,M_{nj_n}):B}
 \\
 \inferrule*[right=$\otimes\mathsf I$]
 {\Gamma\vdash\sum_{i=1}^{m}M_i:A\quad
  \Delta\vdash\sum_{j=1}^{n}N_j:B}
 {\Gamma,\Delta\vdash
  \displaystyle\sum_{i=1}^{m}\sum_{j=1}^{n}\{M_i,N_j\}:A\otimes B}
 \\
 \inferrule*[right=$\otimes\mathsf E$]
 {\Gamma\vdash\sum_{i=1}^{m}M_i:A\otimes B\quad
  \Delta_1,x:A,y:B,\Delta_2\vdash\sum_{j=1}^{n}N_j:C}
 {\Delta_1,\Gamma,\Delta_2\vdash
  \displaystyle\sum_{i=1}^{m}\sum_{j=1}^{n}
  \operatorname{match}(M_i,xy.N_j):C}
 \\
 \inferrule*[right=$\I\mathsf I$]{\;}{()\vdash\star:\I}
 \qquad
 \inferrule*[right=$\I\mathsf E$]
 {\Gamma\vdash\sum_{i=1}^{m}M_i:\I\quad
  \Delta_1,\Delta_2\vdash\sum_{j=1}^{n}N_j:A}
 {\Delta_1,\Gamma,\Delta_2\vdash
  \displaystyle\sum_{i=1}^{m}\sum_{j=1}^{n}
  \operatorname{match}_{\I}(M_i,N_j):A}
 \\
 \inferrule*[right=$\join\mathsf I$]
 {\Gamma\vdash M:A\quad\Gamma\vdash N:A}
 {\Gamma\vdash M\join N:A}.
\end{gather*}
Exchange and general injective renaming are admissible alongside the displayed
weakening rule.  In every join-free summand each free variable occurs at most
once: the premises of $f$, $\otimes\mathsf I$, and the elimination rules have
disjoint contexts.  The two premises of $\join\mathsf I$, by contrast, have
the same context because they describe alternatives rather than simultaneous
uses of their inputs.  This difference is the syntactic content of
affineness.  Nullary generators may be weakened into an arbitrary context.

It is important that the rules themselves contain the joins.  For example,
the conclusion of the $f$-rule is not merely
$f(\sum_jM_{1j},\ldots,\sum_jM_{nj})$ together with later distributivity
equations: it is already the join of all choices.  Thus multilinearity is
built into typing derivations.  Every term can consequently be regarded as a
finite non-empty list of join-free summands, with joins flattened at the top
level.

The equality on terms contains associativity, commutativity, and idempotence
of $\join$, is a congruence for the extended constructors, and contains the
representability equations
\begin{align*}
 \operatorname{match}(\{P,Q\},xy.R)
   &\equiv R[P/x][Q/y],
 &
 \operatorname{match}(S,xy.R[\{x,y\}/u])
   &\equiv R[S/u],\\
 \operatorname{match}_{\I}(\star,R)
   &\equiv R,
 &
 \operatorname{match}_{\I}(S,R[\star/u])
   &\equiv R[S/u].
\end{align*}
The first equation in each pair is a $\beta$-law and the second is an
$\eta$-law.  They express exactly that $\{x,y\}$ and $\star$ are the universal
arrows representing $A\otimes B$ and $\I$.

\subsection{Coherence of derivations}

Because joins may occur in the premises of a rule, a raw proof tree is not
literally determined by its conclusion.  A join may be consumed by one
extended rule instance or first split by $\join\mathsf I$ and distributed
through that rule.  The correct uniqueness statement is therefore the
following.

\begin{proposition}[Coherence of typing derivations]
\label{prop:type-theory-derivation-coherence}
Every derivable term has a unique normal derivation.  Every derivation of that
term is permutation equivalent to its normal derivation; hence a term uniquely
determines a permutation-equivalence class of derivations.
\end{proposition}

Here permutation equivalence is generated by commuting an exposed
$\join\mathsf I$ outwards through any other rule, together with the usual
conversions for injective renaming and weakening.  A derivation is normal when
it is a top-level list of join-free derivations, all constructor premises are
singletons, and unused variables have been moved to a fixed canonical
weakening at the boundary.

The proof is a normalisation argument.  First consider a join-free term in
the context consisting of precisely its free variables.  Its outermost
constructor determines the last rule, affineness determines a unique split of
the context between the premises, and induction determines the subderivations.
There is therefore a unique tight derivation.  Inject this free-variable
context into the displayed ambient context and put the unused variables in a
fixed position; this gives the unique normal join-free derivation.  Induction
on the list of summands gives uniqueness for an arbitrary term.

For an arbitrary derivation, normalise its premises and commute every exposed
join through the last rule in a fixed left-to-right order.  The conclusions of
the rules above are already the corresponding row-major lists of summands, so
this process changes the proof tree but not its conclusion.  Finally, commute
all internal weakenings to the boundary.  The result is the normal derivation
just described.  Any two derivations of the same term therefore meet at this
normal form.

This computational core is mechanised in Agda.  In the formalisation, terms
are non-empty lists of join-free terms, extended binary constructors are
implemented by a canonical row-major product, and a normalisation function
distributes rule applications into a normal derivation.  Uniqueness for
join-free and normal derivations is proved by
\[
 \texttt{JFDerivation-unique}
 \qquad\text{and}\qquad
 \texttt{NFDerivation-unique},
\]
after which \texttt{termsDetermineDerivations} concludes that two derivations
of the same term have equal normal forms.  The current file formalises binary
generators and a common-context presentation; the extra affine step above is
the free-variable support and canonical-weakening argument, which remains to
be added to the mechanisation.

\subsection{Substitution}

Let
\[
 \Gamma\vdash P=\sum_{j=1}^{m}p_j:A,
 \qquad
 \Delta_1,x:A,\Delta_2\vdash Q=\sum_{i=1}^{n}q_i:B,
\]
where $p_j$ and $q_i$ are join-free and the contexts have first been renamed
apart.  Since a join-free affine term contains at most one occurrence of $x$,
ordinary capture-avoiding substitution $q_i\{p_j/x\}$ is defined without
duplicating $p_j$.  If $q_i$ does not use $x$, the substitution simply leaves
$q_i$ unchanged.  Define substitution of general terms by
\begin{equation}
\label{eq:type-theory-substitution}
 Q[P/x]
 \defeq
 \sum_{i=1}^{n}\sum_{j=1}^{m}q_i\{p_j/x\}.
\end{equation}
On derivations, normalise both arguments, substitute every pair of join-free
summands, and join the resulting singleton derivations in the displayed
order.  Proposition~\ref{prop:type-theory-derivation-coherence} makes the
result independent, up to permutation equivalence, of the input derivations.
Notice that if $q_i$ discards $x$,~\eqref{eq:type-theory-substitution} contains
one copy of $q_i$ for every $p_j$; idempotence makes these copies equal to a
single $q_i$, as required for affine weakening.

The usual substitution lemma is then admissible:
\[
 \inferrule
 {\Gamma\vdash P:A\quad
  \Delta_1,x:A,\Delta_2\vdash Q:B}
 {\Delta_1,\Gamma,\Delta_2\vdash Q[P/x]:B}.
\]
Writing $M\equiv_{\join}N$ when two flattened terms have the same set of
summands, substitution satisfies, whenever the expressions are well typed,
\begin{align*}
 x[P/x]&=P,
 &Q[x/x]&=Q,\\
 (Q\join R)[P/x]&=Q[P/x]\join R[P/x],
 &Q[(P\join R)/x]&\equiv_{\join}Q[P/x]\join Q[R/x],\\
 (Q[P/x])[N/y]&\equiv_{\join}Q[P[N/y]/x],
 &Q[P/x][N/y]&\equiv_{\join}Q[N/y][P/x].
\end{align*}
The last two equations are, respectively, associativity and interchange.  The
freshness conditions are the evident ones: in the associativity equation $y$
belongs to the context of $P$ and is fresh for $Q$; in the interchange
equation $x$ and $y$ are distinct inputs of $Q$ and neither substituted term
contains variables from the other's context.

These equations are not merely convenient metatheory.  The two unit equations
give identities for placed multicategorical composition; associativity and
interchange are its remaining axioms; and the two join equations say that
composition is a semilattice homomorphism in each argument.  Substitution also
commutes with injective reindexing: renaming and adding unused inputs before
substitution gives the same result as applying the induced injection to the
composite context afterwards.  These are precisely the operations and axioms
required of an affine $\catname{SLat}$-multicategory.  Finally, the four
$\beta$--$\eta$ equations above supply the chosen representations for
$\otimes$ and $\I$.

Term equality must be stable under substitution for composition to descend to
equivalence classes.  In the affine calculus this follows by induction on the
generation of equality.  The semilattice cases use the two join laws.  In a
congruence case there is at most one path from the root of a join-free term to
the substituted variable, so the equality may be transported along this
one-hole context.  The $\beta$--$\eta$ cases reduce to associativity of
join-free substitution, and injective reindexing commutes with substitution as
described above.  This is also the precise point at which contraction would
break the proof by creating more than one such path.

\subsection{Expected universal property}

The preceding construction supplies all the data of an affine representable
multicategory enriched in $\catname{SLat}$.  Its objects are the generated types, its
hom-semilattices are equivalence classes of judgments, identities are
variables, composition is substitution, and the structural action is
injective renaming and weakening.  Pairing and the unit term give the chosen
representations.  What remains is to assemble these observations into the
universal property and to mechanise the affine structural part of the
coherence proof.  We therefore end with the statement as a conjecture rather
than a result of this thesis.

\paragraph{Conjecture (initiality).}
For every multigraph $\mathcal G$, the generated types and equivalence classes
of derivable terms form the free representable affine
$\catname{SLat}$-multicategory on $\mathcal G$.  Equivalently, this construction
is the left adjoint in
\[
 \catname{Multigraph}
 \;\underset{U}{\stackrel{F}{\rightleftarrows}}\;
 \catname{SLat}\text{-}\catname{AffRepMultCat}.
\]

The expected proof interprets a join-free term in any representable affine
$\catname{SLat}$-multicategory by induction on its syntax, using
multicategorical composition for generators, the universal arrows for pairing
and matching, and the injection action for unused variables.  A general term
is interpreted by joining the interpretations of its summands.  Coherence
shows that this interpretation does not depend on a typing derivation, and the
substitution lemma shows that it preserves composition.  The semilattice and
$\beta$--$\eta$ equations make it well defined on equivalence classes.
Structural induction then gives uniqueness.

This conjecture captures the part of the original project that appears to
survive its cartesian obstruction.  It does not recover arbitrary e-graph
sharing or cycles: those require a separate account of correlated choice,
contraction, and guarded recursion.  Its value is instead to isolate a sound
type-theoretic core and, in the symmetric restriction, to provide a compact
term language for semilattice-enriched string diagrams.
